\documentclass{amsart}
% For dealing with references we use the comment environment
\usepackage{verbatim}
\newenvironment{reference}{\comment}{\endcomment}
%\newenvironment{reference}{}{}
\newenvironment{slogan}{\comment}{\endcomment}
\newenvironment{history}{\comment}{\endcomment}

% For commutative diagrams we use Xy-pic
\usepackage[all]{xy}

% We use 2cell for 2-commutative diagrams.
\xyoption{2cell}
\UseAllTwocells

% We use multicol for the list of chapters between chapters
\usepackage{multicol}

% This is generally recommended for better output
\usepackage{lmodern}
\usepackage[T1]{fontenc}

% For cross-file-references

% Package for hypertext links:
\usepackage{hyperref}

% For any local file, say "hello.tex" you want to link to please

% Theorem environments.
%
\theoremstyle{plain}
\newtheorem{theorem}[subsection]{Theorem}
\newtheorem{proposition}[subsection]{Proposition}
\newtheorem{lemma}[subsection]{Lemma}

\theoremstyle{definition}
\newtheorem{definition}[subsection]{Definition}
\newtheorem{example}[subsection]{Example}
\newtheorem{exercise}[subsection]{Exercise}
\newtheorem{situation}[subsection]{Situation}

\theoremstyle{remark}
\newtheorem{remark}[subsection]{Remark}
\newtheorem{remarks}[subsection]{Remarks}

\numberwithin{equation}{subsection}

% Macros
%
\def\lim{\mathop{\mathrm{lim}}\nolimits}
\def\colim{\mathop{\mathrm{colim}}\nolimits}
\def\Spec{\mathop{\mathrm{Spec}}}
\def\Hom{\mathop{\mathrm{Hom}}\nolimits}
\def\Ext{\mathop{\mathrm{Ext}}\nolimits}
\def\SheafHom{\mathop{\mathcal{H}\!\mathit{om}}\nolimits}
\def\SheafExt{\mathop{\mathcal{E}\!\mathit{xt}}\nolimits}
\def\Sch{\mathit{Sch}}
\def\Mor{\mathop{\mathrm{Mor}}\nolimits}
\def\Ob{\mathop{\mathrm{Ob}}\nolimits}
\def\Sh{\mathop{\mathit{Sh}}\nolimits}
\def\NL{\mathop{N\!L}\nolimits}
\def\CH{\mathop{\mathrm{CH}}\nolimits}
\def\proetale{{pro\text{-}\acute{e}tale}}
\def\etale{{\acute{e}tale}}
\def\QCoh{\mathit{QCoh}}
\def\Ker{\mathop{\mathrm{Ker}}}
\def\Im{\mathop{\mathrm{Im}}}
\def\Coker{\mathop{\mathrm{Coker}}}
\def\Coim{\mathop{\mathrm{Coim}}}

% Boxtimes
%
\DeclareMathSymbol{\boxtimes}{\mathbin}{AMSa}{"02}

%
% Macros for moduli stacks/spaces
%
\def\QCohstack{\mathcal{QC}\!\mathit{oh}}
\def\Cohstack{\mathcal{C}\!\mathit{oh}}
\def\Spacesstack{\mathcal{S}\!\mathit{paces}}
\def\Quotfunctor{\mathrm{Quot}}
\def\Hilbfunctor{\mathrm{Hilb}}
\def\Curvesstack{\mathcal{C}\!\mathit{urves}}
\def\Polarizedstack{\mathcal{P}\!\mathit{olarized}}
\def\Complexesstack{\mathcal{C}\!\mathit{omplexes}}
% \Pic is the operator that assigns to X its picard group, usage \Pic(X)
% \Picardstack_{X/B} denotes the Picard stack of X over B
% \Picardfunctor_{X/B} denotes the Picard functor of X over B
\def\Pic{\mathop{\mathrm{Pic}}\nolimits}
\def\Picardstack{\mathcal{P}\!\mathit{ic}}
\def\Picardfunctor{\mathrm{Pic}}
\def\Deformationcategory{\mathcal{D}\!\mathit{ef}}

\setlength{\emergencystretch}{3em}
\begin{document}
\title[Stacks Project, Tag 0GA1]{315 --- Uniform annihilation of high Tor against every ideal-power quotient}
\maketitle
\noindent Copyright (C) 2005--2025 Johan de Jong. Stacks Project authors. Source: \href{https://stacks.math.columbia.edu/tag/0GA1}{Tag 0GA1}.

\noindent Editorial difficulty note (not author text). Difficulty H2: The proof represents derived tensor with ideal powers as cohomology on the blowup. A bounded spectral sequence assembles sheafwise annihilators into a uniform global exponent, independent of the module, homological degree and power of the ideal. The uniformity and filtration control remain substantive after the local blowup estimates..

\noindent Editorial provenance note (not author text). History: excerpt from revision \texttt{a0444\allowbreak 6e57e\allowbreak c1fbc\allowbreak 252a8\allowbreak 71afc\allowbreak ec775\allowbreak 2fb28\allowbreak 07b14}, local-cohomology.tex, lines 4451--4494. Reference presentation was adapted; original-author context and core support blocks were retained or added. The mathematical wording is unchanged. Licensed under GNU FDL 1.2 or later; no invariant sections or cover texts. See ../licenses/COPYING. Context blocks reproduce author definitions and supporting results; they are not additional counted problems.

\section{A bit of uniformity, II}
\label{section-uniform-vanishing-bis}

\noindent
Let $I$ be an ideal of a Noetherian ring $A$. Let $M$ be a finite
$A$-module. Let $i > 0$. By More on Algebra, Lemma
\href{https://stacks.math.columbia.edu/tag/0911}{lemma tor strictly pro zero (Tag 0911)}
there exists a $c = c(A, I, M, i)$ such that
$\text{Tor}^A_i(M, A/I^n) \to \text{Tor}^A_i(M, A/I^{n - c})$
is zero for all $n \geq c$. In this section, we discuss
some results which show that one sometimes can choose
a constant $c$ which works for all $A$-modules $M$ simultaneously
(and for a range of indices $i$).
This material is related to uniform Artin-Rees as discussed in
\href{https://stacks.math.columbia.edu/bibliography}{Bibliography: Huneke-uniform} and \href{https://stacks.math.columbia.edu/bibliography}{Bibliography: AHS}.

\medskip\noindent
In Remark \href{https://stacks.math.columbia.edu/tag/0GA5}{remark strict pro isomorphism (Tag 0GA5)} we will apply this to show that
various pro-systems related to derived completion are (or are not)
strictly pro-isomorphic.

\medskip\noindent
The following lemma can be significantly strengthened.



\begin{lemma}
\label{lemma-annihilates-tors}
In the situation above, let $t$ be an upper bound on the number of
generators for $I$. There exists an integer $c = c(A, I) \geq 0$
such that for any $A$-module $M$ the tor modules
$\text{Tor}_i^A(M, A/I^q)$ are annihilated by $I^c$ for $i > t$
and all $q \geq 0$.
\end{lemma}

\begin{proof}
Let $q(A, I)$ be as above. For $q \geq q(A, I)$ we have
$$
R\Gamma(X, Lp^*\widetilde{M}(q)) = M \otimes_A^\mathbf{L} I^q
$$
by Lemma \href{https://stacks.math.columbia.edu/tag/0G9Z}{lemma compute tor Iq (Tag 0G9Z)}.
We have a bounded and convergent spectral sequence
$$
H^a(X, H^b(Lp^*\widetilde{M}(q))) \Rightarrow
\text{Tor}_{-a - b}^A(M, I^q)
$$
by Derived Categories of Schemes, Lemma \href{https://stacks.math.columbia.edu/tag/0G9Q}{lemma spectral sequence (Tag 0G9Q)}.
Let $d$ be an integer as in Cohomology of Schemes, Lemma
\href{https://stacks.math.columbia.edu/tag/071L}{lemma vanishing nr affines quasi separated (Tag 071L)}
(actually we can take $d = t$, see
Cohomology of Schemes, Lemma \href{https://stacks.math.columbia.edu/tag/01XI}{lemma vanishing nr affines (Tag 01XI)}).
Then we see that $H^{-i}(X, Lp^*\widetilde{M}(q)) = \text{Tor}_i^A(M, I^q)$
has a finite filtration with at most $d$ steps whose graded are
subquotients of the modules
$$
H^a(X, H^{- i - a}(Lp^*\widetilde{M})(q)),\quad
a = 0, 1, \ldots, d - 1
$$
If $i \geq t$ then all of these modules are annihilated
by $I^c$ where $c = c(A, I)$ is as in Lemma \href{https://stacks.math.columbia.edu/tag/0GA0}{lemma annihilates (Tag 0GA0)}
because the cohomology sheaves $H^{- i - a}(Lp^*\widetilde{M})$
are all annihilated by $I^c$ by the lemma. Hence we see that
$\text{Tor}_i^A(M, I^q)$ is annihilated by $I^{dc}$ for
$q \geq q(A, I)$ and $i \geq t$. Using the short exact sequence
$0 \to I^q \to A \to A/I^q \to 0$ we find that
$\text{Tor}_i(M, A/I^q)$ is annihilated by $I^{dc}$ for $q \geq q(A, I)$
and $i > t$. We conclude that $I^m$ with $m = \max(dc, q(A, I) - 1)$
annihilates $\text{Tor}_i^A(M, A/I^q)$ for all $q \geq 0$
and $i > t$ as desired.
\end{proof}
\end{document}
